\documentclass[aspectratio=169,11pt]{beamer}
\usepackage{fontspec}
\usepackage[brazil]{babel}
\usetheme{Madrid}
\usecolortheme{whale}
\usefonttheme{professionalfonts}
\usepackage{pgfplots}\pgfplotsset{compat=1.18}
\pgfplotsset{/pgf/number format/assume math mode=true} % dígitos dos eixos em modo texto (evita nullfont no Beamer+XeTeX)
\usepackage{tikz}\usetikzlibrary{automata,positioning,arrows.meta}
\usepackage{booktabs}
\usepackage{listings}
\usepackage{unicode-math}
\setmathfont{latinmodern-math.otf} % fonte de matemática OpenType: corrige dígitos em math no Beamer+XeTeX

\lstset{basicstyle=\ttfamily\scriptsize,breaklines=true,columns=fullflexible}

\title[Navegação Autônoma --- Khepera IV]{Sistema de Navegação Autônoma para o Robô Móvel Khepera~IV}
\subtitle{Projeto Final --- FTL094 Engenharia de Software de Tempo Real}
\author[Equipe FTL094]{\small Gabriel Henrique de Souza Nazaré \and Luan Alexandre Ramos Barreto \and Matheus Rocha Canto \and Matheus Serrão Uchôa \and Ricardo Mendonça Braz \and Yan Fernandes da Silva}
\institute[UFAM/CETELI]{PPGEE --- CETELI --- Universidade Federal do Amazonas}
\date{Manaus--AM, 2026}

\begin{document}

\begin{frame}\titlepage\end{frame}

\begin{frame}{Agenda}\tableofcontents\end{frame}

% ---------------------------------------------------------------
\section{Contexto e Objetivos}
\begin{frame}{O problema}
\begin{columns}
\column{0.58\textwidth}
\begin{itemize}
  \item \textbf{Encomenda:} software de um \emph{robô móvel autônomo}.
  \item Levar o robô de um ponto \textbf{A} a um ponto \textbf{B}\\ \emph{desviando de todos os obstáculos}.
  \item Plataforma: \textbf{Khepera~IV} (CETELI/UFAM).
  \item Desafio: sensoriamento \textbf{curto} ($\approx$2--20\,cm) e \textbf{sem mapa} $\Rightarrow$ navegação reativa.
\end{itemize}
\column{0.40\textwidth}
\begin{block}{Objetivo}
Atingir o alvo em cenário com 5 obstáculos, \textbf{sem colisões}, em laço de \textbf{tempo real} (16~ms).
\end{block}
\end{columns}
\end{frame}

% ---------------------------------------------------------------
\section{Requisitos}
\begin{frame}{Requisitos (resumo)}
\begin{columns}[T]
\column{0.5\textwidth}
\textbf{Funcionais}
\begin{itemize}\small
  \item RF-01 ir de A a B
  \item RF-02 desviar de obstáculos
  \item RF-03 detectar chegada e parar
  \item RF-06 escapar de mínimo local
  \item RF-07 telemetria
\end{itemize}
\column{0.5\textwidth}
\textbf{Não funcionais / Tempo real}
\begin{itemize}\small
  \item RNF-01 sem colisões
  \item RNF-04 portável (sim$\to$real)
  \item RT-01 período fixo de 16~ms
  \item RT-02 \emph{deadline} a cada ciclo
\end{itemize}
\end{columns}
\end{frame}

% ---------------------------------------------------------------
\section{Arquitetura}
\begin{frame}{Arquitetura: \emph{Sentir--Decidir--Atuar} + FSM}
\centering
\begin{tikzpicture}[comp/.style={draw,rounded corners,minimum height=0.9cm,minimum width=2.1cm,align=center,font=\scriptsize,fill=blue!8},
  io/.style={draw,minimum height=0.9cm,minimum width=1.7cm,align=center,font=\scriptsize,fill=black!8},>={Stealth},node distance=0.6cm]
  \node[io] (s) {Sensores};
  \node[comp,right=of s] (p) {Percepção};
  \node[comp,right=of p] (f) {FSM};
  \node[comp,right=of f] (c) {Controle};
  \node[io,right=of c] (m) {Motores};
  \draw[->](s)--(p);\draw[->](p)--(f);\draw[->](f)--(c);\draw[->](c)--(m);
\end{tikzpicture}
\vspace{0.4cm}
\begin{tikzpicture}[->,>=Stealth,node distance=2.9cm,on grid,
  every state/.style={draw,rounded corners,minimum width=1.9cm,align=center,font=\scriptsize}]
  \node[state] (N) {\textsc{Navegar}};
  \node[state] (G) [right=of N] {\textsc{Girar}};
  \node[state] (C) [right=of G] {\textsc{Contornar}};
  \path (N) edge[bend left] node[above,font=\tiny]{frente$>$trig} (G)
        (G) edge[bend left] node[above,font=\tiny]{frente$<$clear} (C)
        (C) edge[bend left] node[below,font=\tiny]{nova parede} (G)
        (C) edge[bend left=55] node[below,font=\tiny]{passou $d_p$} (N);
\end{tikzpicture}
\end{frame}

% ---------------------------------------------------------------
\section{Implementação}
\begin{frame}[fragile]{Implementação}
\begin{columns}[T]
\column{0.5\textwidth}
\begin{itemize}\small
  \item Simulador \textbf{Webots R2025a}.
  \item Controlador em \textbf{C} (toolchain do Webots).
  \item Localização: GPS+IMU (sim.); odometria (real).
  \item Lei de rumo P + repulsão reativa por IR.
\end{itemize}
\column{0.5\textwidth}
\begin{block}{Lei de controle}
\[\omega = K_\theta\,e_\theta + K_a\,(b_R-b_L)\]
\[\dot\phi_L=v-\omega,\quad \dot\phi_R=v+\omega\]
\end{block}
\end{columns}
\end{frame}

% ---------------------------------------------------------------
\section{Engenharia e Depuração}
\begin{frame}[fragile]{Duas falhas reais (e suas lições)}
\begin{block}{1. Mínimo local --- robô preso (telemetria)}
\begin{lstlisting}
pos=(-0.54,-0.54) dist=1.90 err=+0.00 IR_F=280 rodas=-3.0/+3.0
pos=(-0.54,-0.54) dist=1.90 err=-0.01 IR_F=279 rodas=+3.0/-3.0
\end{lstlisting}
\footnotesize Posição congelada, rodas oscilando $\Rightarrow$ campo potencial travado. \textbf{Cura:} FSM com contorno comprometido (Bug).
\end{block}
\begin{block}{2. Calibração: folha de dados $\neq$ realidade}
\footnotesize Base do IR prevista $\approx 29$, \textbf{medida $\approx 110$} (reflexão do piso). \textbf{Lição:} medir, nunca supor.
\end{block}
\end{frame}

% ---------------------------------------------------------------
\section{Resultados}
\begin{frame}{Resultado: trajetória livre de colisões}
\begin{columns}
\column{0.6\textwidth}
\centering
\begin{tikzpicture}
\begin{axis}[axis equal image,width=\textwidth,xmin=-1.6,xmax=1.6,ymin=-1.6,ymax=1.6,
  xlabel={$x$ (m)},ylabel={$y$ (m)},tick label style={font=\tiny},label style={font=\scriptsize},grid=both,grid style={gray!20}]
  \draw[thick,gray] (axis cs:-1.5,-1.5) rectangle (axis cs:1.5,1.5);
  \fill[brown!55](axis cs:-0.6,-0.6)rectangle(axis cs:-0.4,-0.4);
  \fill[brown!55](axis cs:-0.125,-0.125)rectangle(axis cs:0.125,0.125);
  \fill[brown!55](axis cs:0.4,0.4)rectangle(axis cs:0.6,0.6);
  \fill[brown!55](axis cs:0.325,-0.375)rectangle(axis cs:0.475,-0.225);
  \fill[brown!55](axis cs:-0.375,0.325)rectangle(axis cs:-0.225,0.475);
  \addplot[blue,thick,mark=*,mark size=0.7pt] coordinates {
  (-1.00,-1.00)(-0.85,-0.85)(-0.70,-0.70)(-0.67,-0.62)(-0.79,-0.53)(-0.92,-0.44)
  (-0.91,-0.32)(-0.80,-0.18)(-0.63,-0.05)(-0.45,0.07)(-0.27,0.18)(-0.08,0.29)
  (0.09,0.41)(0.27,0.53)(0.30,0.66)(0.29,0.84)(0.34,0.95)(0.51,1.01)(0.71,1.03)(0.94,1.03)};
  \addplot[only marks,mark=triangle*,green!50!black,mark size=3pt]coordinates{(-1,-1)};
  \addplot[only marks,mark=square*,red,mark size=3pt]coordinates{(1,1)};
\end{axis}
\end{tikzpicture}
\column{0.4\textwidth}
\small
\begin{tabular}{@{}ll@{}}
\toprule
Métrica & Valor \\
\midrule
Início & $(-1,-1)$ \\
Chegada & $(0.94,1.03)$ \\
Dist.\ inicial & $2.83$~m \\
Tempo & $19.3$~s \\
Colisões & \textbf{0} \\
Obstáculos & 5 \\
\bottomrule
\end{tabular}
\end{columns}
\end{frame}

\begin{frame}{Resultado: convergência ao alvo}
\centering
\begin{tikzpicture}
\begin{axis}[width=0.72\textwidth,height=5cm,xlabel={tempo $t$ (s)},ylabel={distância (m)},
  xmin=0,xmax=20,ymin=0,ymax=3,grid=both,grid style={gray!20},tick label style={font=\scriptsize},label style={font=\small}]
  \addplot[blue,thick,mark=*,mark size=0.8pt] coordinates {
  (0,2.83)(1,2.62)(2,2.40)(3,2.33)(4,2.36)(5.1,2.40)(6.1,2.33)(7.1,2.15)(8.1,1.94)
  (9.1,1.73)(10.1,1.51)(11.1,1.30)(12.1,1.08)(13.1,0.86)(14.1,0.77)(15.1,0.73)
  (16.1,0.66)(17.2,0.49)(18.2,0.29)(19.2,0.09)(19.3,0.0)};
  \addplot[red,dashed]coordinates{(0,0.07)(20,0.07)};
\end{axis}
\end{tikzpicture}\\
\footnotesize Redução quase monotônica; o leve aumento entre 3 e 6~s é o custo do contorno da 1ª caixa.
\end{frame}

% ---------------------------------------------------------------
\section{Verificação e Validação}
\begin{frame}{Verificação e validação}
\begin{columns}[T]
\column{0.52\textwidth}
\textbf{Casos de teste (Webots)} --- 6/6 \textbf{PASSOU}
\begin{itemize}\scriptsize
  \item CT-01 caminho livre \hfill \textbf{ok}
  \item CT-02 desvio de obstáculo \hfill \textbf{ok}
  \item CT-03 missão no \emph{slalom} \hfill \textbf{ok}
  \item CT-04 escape de mínimo local \hfill \textbf{ok}
  \item CT-05 período de 16~ms \hfill \textbf{ok}
  \item CT-06 telemetria e calibração \hfill \textbf{ok}
\end{itemize}
\scriptsize 3 defeitos diagnosticados por telemetria e corrigidos.
\column{0.48\textwidth}
\begin{block}{Bateria de 1000 mundos (OpenGL)}
\scriptsize O simulador \textbf{reusa o mesmo \texttt{controller\_core.h}} do robô físico. 1000 cenários de dificuldade progressiva (até 33 caixas), \emph{seed} fixa. Desfecho: \textsc{Arrived} / \textsc{Collision} / \textsc{Timeout}.
\end{block}
\scriptsize Chega na maioria; falhas concentram-se nos mundos densos (limite teórico do Bug2 reativo).
\end{columns}
\end{frame}

% ---------------------------------------------------------------
\section{Evolução: Robô Físico}
\begin{frame}{Evolução para o Khepera~IV físico}
\begin{columns}[T]
\column{0.54\textwidth}
\begin{itemize}\small
  \item \textbf{Odometria} (encoders: 147,4 pulsos/mm; base 105,4~mm) no lugar do GPS.
  \item \textbf{Bug2} canônico: segue a linha $A\!\to\!B$ e larga a parede ao \emph{recruzá-la mais perto de $B$}.
  \item \textbf{Ponto de controle virtual (ponto-P):} volta ao eixo $A\!\to\!B$ após o contorno (PI + anti-windup).
  \item \textbf{Camada deliberativa:} mapa de ocupação + planejamento global \textbf{$A^\star$}.
\end{itemize}
\column{0.46\textwidth}
\begin{block}{Interface web (Flask)}
\scriptsize Conecta por USB/Wi-Fi; \emph{deploy}, recompilação e ajuste de parâmetros no robô; telemetria; teleop; mapeamento; navegação por $A^\star$; painel do simulador OpenGL.
\end{block}
\scriptsize Laço embarcado periódico ($\approx$40~ms) --- tarefa de tempo real.
\end{columns}
\end{frame}

% ---------------------------------------------------------------
\section{Conclusão}
\begin{frame}{Conclusão e trabalhos futuros}
\begin{block}{Conclusão}
Arquitetura híbrida (FSM) + calibração empírica $\Rightarrow$ alvo atingido sem colisões em 19,3~s, respeitando o laço de tempo real. Depuração guiada por telemetria foi decisiva. Validado em campanha de testes e em \textbf{1000 cenários}; transferido para o \textbf{robô físico}.
\end{block}
\begin{block}{Já concretizado (era ``trabalho futuro'')}
Odometria \& Bug2 (robô real) $\bullet$ rastreamento por ponto-P $\bullet$ mapa + $A^\star$ $\bullet$ interface web.
\end{block}
\begin{block}{Trabalhos futuros remanescentes}
Ultrassom (antecipação) $\bullet$ VFH/DWA $\bullet$ localização persistente $\bullet$ análise formal de escalonabilidade (RMS).
\end{block}
\end{frame}

\begin{frame}[plain]
\centering
{\Huge Obrigado!}\\[8pt]
{\large Perguntas?}\\[16pt]
{\footnotesize Equipe FTL094 --- PPGEE/CETELI/UFAM (2026/1)}
\end{frame}

\end{document}
